% Source: Wald GR, physical PDF p. 118
%         (printed p. 105).
% Coverage: Figure 5.6.
% Figures/tables/footnotes: Figure 5.6; no tables or footnotes.
% Status: source-reviewed and compile-clean.
% Uncertainties: none.
\begingroup
\renewcommand{\thefigure}{5.6}
\begin{figure}[htbp]
  \centering
  \begin{tikzpicture}[>=Stealth, line cap=round, line join=round, scale=.83]
    \fill[gray!18] (-3.55,0) rectangle (3.55,.43);
    \draw[thick] (-3.55,.43) -- (3.55,.43);
    \foreach \x in {-2.70,-1.35,1.35,2.70} {
      \draw[thick] (\x,.43) .. controls (\x-.11,1.65) and (\x+.08,2.75) .. (\x+.28,4.10);
    }
    \fill (0,1.88) circle (1.5pt);
    \draw[dashed, thick] (0,1.88) -- (-2.12,.43);
    \draw[dashed, thick] (0,1.88) -- (2.12,.43);
    \draw[thick] (0,.43) -- (0,1.88);
    \node[fill=white,inner sep=1.5pt] at (0,.18) {singularity};
  \end{tikzpicture}
  \caption{The causal structure of the Robertson--Walker solutions near the big-bang
  singularity. Particle horizons exist since an observer cannot ``see'' all other
  isotropic observers in the universe.}
  \label{fig:5.6}
\end{figure}
\endgroup
